\begin{tikzpicture}[
        scale=1.5,
        >=Stealth,
        % 随机游走路径样式
        rw-path/.style={thick, blue!70!black, mark=*, mark size=1.5pt},
        % 管道填充样式
        tube-fill/.style={gray!10},
        % 标注样式
        label-font/.style={font=\small}
    ]

        % 1. 绘制背景管道填充 (r_0 到 r_0 + delta)
        \fill[tube-fill] (0, -1.5) rectangle (4, 0.5);

        % 2. 绘制辅助虚线
        \draw[dashed, black!60] (0, -2) -- (4.5, -2) node[right] {$-\Delta$};
        % \draw[dashed, black!60] (0, 0) -- (4.5, 0) node[right] {$y = 0$};
        \draw[dashed, black!60] (4, -2.5) -- (4, 1) node[above] at (4, 1) {$L$};

        % 3. 绘制管道实线 (r_0 到 r_0 + delta)
        \draw[thick, black] (0, -1.5) -- (4.5, -1.5) node[right] {$r_0$};
        \draw[thick, black] (0, 0.5) -- (4.5, 0.5) node[right] {$r_0 + \Delta$};

        % 4. 绘制坐标轴 (Y轴实线)
        \draw[->, thick] (0, -2.5) -- (0, 1.2) node[above] {$Y_n$};
        \draw[->, thick, dashed] (0, 0) -- (5, 0) node[right] {$n$};

        % 5. 绘制随机游走 (起点在 0,0，始终保持在 [-1.5, 0.5] 之间)
        % 选取了一些看起来随机的点
        \draw[rw-path] plot coordinates {
            (0, 0) 
            (0.8, -0.6) 
            (1.5, 0.2) 
            (2.2, -1.1) 
            (3.1, -0.4) 
            (4, -0.8)
        };

        % 6. 标注 Delta (使用大括号)
        \draw[decorate, decoration={brace, amplitude=5pt}] 
            (-0.1, -1.5) -- (-0.1, 0.5) node[midway, left=6pt] {$\Delta$};

        % 标注特定的点
        % \node[below left] at (0,0) {$0$};

    \end{tikzpicture}